\section{Diseño del Proyecto}

La fase de diseño del proyecto GymFlow tiene como objetivo definir la estructura conceptual y visual del sistema antes de su implementación. En esta etapa se establecen los lineamientos generales de la interfaz de usuario, la navegación dentro de la aplicación y la organización de los datos que serán utilizados por el sistema.

En la versión actual del proyecto ya se cuenta con los mockups principales de la aplicación. Los diagramas formales de navegación y del modelo de datos todavía se encuentran en elaboración, por lo que en esta sección se presenta el diseño funcional definido hasta el momento y la estructura conceptual que guiará la implementación.

\subsection{Diseño de interfaces}

El diseño de interfaces de \textbf{GymFlow} se define con un enfoque funcional y directo, orientado al uso durante la sesión real de entrenamiento. En esta etapa, la prioridad no es la complejidad visual, sino la facilidad de uso para que el usuario pueda continuar su rutina cuando un ejercicio o máquina no se encuentre disponible.

Desde el enfoque visual, la propuesta adopta una interfaz móvil moderna con predominio del color azul, tarjetas limpias, jerarquía tipográfica clara y una navegación enfocada en reducir pasos innecesarios. Este criterio permite mantener una experiencia consistente con el alcance del producto mínimo viable (MVP), pero con una presentación más cercana a un prototipo de alta fidelidad.

La pantalla de \textit{login} permite ingresar correo electrónico y contraseña, y dispone del botón \textit{Entrar} como acción principal de acceso. Una vez autenticado, el usuario llega a la pantalla \textit{home}, que constituye la pantalla principal del sistema y presenta una lista simple de grupos musculares (por ejemplo, pecho, espalda y pierna) para iniciar la consulta.

Al seleccionar un músculo, la pantalla de ejercicios por músculo muestra una lista de ejercicios asociados, como \textit{press banca}, \textit{aperturas} o \textit{flexiones}. Posteriormente, al seleccionar un ejercicio, la pantalla de variantes presenta alternativas para continuar el entrenamiento (por ejemplo, \textit{press banca} con opciones como \textit{press con mancuernas} o \textit{flexiones}). Esta pantalla representa el componente clave de la aplicación, ya que responde directamente al problema principal del proyecto: mantener la continuidad de la rutina ante la ocupación de equipos.

De forma complementaria, la pantalla de crear rutina permite seleccionar ejercicios y guardar una rutina mediante un nombre. La pantalla de entrenamiento muestra la lista de ejercicios de la rutina activa y permite marcarlos como realizados. Finalmente, la pantalla de historial presenta un registro básico de rutinas completadas (por ejemplo, \textit{Rutina 1 - hoy} y \textit{Rutina 2 - ayer}).

En conjunto, estas siete pantallas conforman una secuencia clara, consistente y suficiente para la versión inicial del sistema. Los layouts de referencia correspondientes a cada interfaz se presentan como \textbf{Figura \ref{fig:layout-login}}, \textbf{Figura \ref{fig:layout-home}}, \textbf{Figura \ref{fig:layout-ejercicios}}, \textbf{Figura \ref{fig:layout-variantes}}, \textbf{Figura \ref{fig:layout-crear-rutina}}, \textbf{Figura \ref{fig:layout-entrenamiento}} y \textbf{Figura \ref{fig:layout-historial}}, con el fin de respaldar visualmente las decisiones de diseño adoptadas.

\begin{figure}[H]
    \centering
    \includegraphics[width=0.35\textwidth]{login.png}
    \caption{Layout de referencia de la pantalla de inicio de sesión (Login).}
    \label{fig:layout-login}
\end{figure}

\begin{figure}[H]
    \centering
    \includegraphics[width=0.35\textwidth]{home.png}
    \caption{Layout de referencia de la pantalla principal de músculos (Home).}
    \label{fig:layout-home}
\end{figure}

\begin{figure}[H]
    \centering
    \includegraphics[width=0.35\textwidth]{ejercicios.png}
    \caption{Layout de referencia de la pantalla de ejercicios por músculo.}
    \label{fig:layout-ejercicios}
\end{figure}

\begin{figure}[H]
    \centering
    \includegraphics[width=0.35\textwidth]{variantes.png}
    \caption{Layout de referencia de la pantalla de variantes de ejercicio.}
    \label{fig:layout-variantes}
\end{figure}

\begin{figure}[H]
    \centering
    \includegraphics[width=0.35\textwidth]{crear-rutina.png}
    \caption{Layout de referencia de la pantalla de creación de rutina.}
    \label{fig:layout-crear-rutina}
\end{figure}

\begin{figure}[H]
    \centering
    \includegraphics[width=0.35\textwidth]{entrenamiento.png}
    \caption{Layout de referencia de la pantalla de entrenamiento.}
    \label{fig:layout-entrenamiento}
\end{figure}

\begin{figure}[H]
    \centering
    \includegraphics[width=0.35\textwidth]{historial.png}
    \caption{Layout de referencia de la pantalla de historial.}
    \label{fig:layout-historial}
\end{figure}

\subsection{Flujo de navegación del sistema}

El flujo de navegación del sistema describe la manera en que los usuarios interactúan con la aplicación y cómo se conectan las diferentes pantallas. La navegación se implementa con \textit{Expo Router} mediante una pila principal (\textit{Stack}) y una barra de pestañas inferior (\textit{Tabs}). El acceso a las distintas zonas está controlado por dos guardas: una de \textbf{autenticación}, que redirige a los usuarios no autenticados al inicio de sesión, y otra de \textbf{onboarding}, que obliga a completar la configuración inicial del perfil antes de ingresar a la aplicación principal.

\subsubsection{Estructura de navegación}

La aplicación se organiza en cuatro grupos de pantallas, descritos en la Tabla \ref{tab:nav-grupos}.

\begin{table}[H]
\centering
\rowcolors{2}{gymblueSoft}{gymblueSoftAlt}
\arrayrulecolor{gymblue}
\begin{tabularx}{0.98\textwidth}{L{3.6cm} Y}
\toprule
\gymtableheaderrow
\gymth{Grupo / Pantalla} & \gymth{Función dentro de la navegación} \\
\midrule
Acceso (\textit{auth}) & Inicio de sesión y registro, incluido el acceso con Google. \\
Onboarding & Configuración inicial guiada del perfil (objetivo, nivel, días y equipamiento) y elección de la rutina inicial. \\
Pestañas principales & Navegación central de la aplicación: \textit{Hoy}, \textit{Rutinas}, \textit{Historial} y \textit{Músculos}. \\
Pantallas de detalle & Ejercicios por músculo, variantes, creación y edición de rutina, y configuración del perfil. \\
\bottomrule
\end{tabularx}
\caption{Grupos de navegación de la aplicación GymFlow}
\label{tab:nav-grupos}
\end{table}

\subsubsection{Recorrido principal del usuario}

En términos generales, el flujo de interacción sigue la siguiente estructura:

\begin{enumerate}
    \item El usuario inicia sesión con correo y contraseña, o con su cuenta de Google.
    \item Si es la primera vez, completa el \textit{onboarding} y obtiene una rutina inicial activa.
    \item En la pestaña \textit{Hoy} visualiza la sesión del día dentro de su calendario semanal.
    \item Durante la sesión, marca series y ejercicios, registra peso, repeticiones y RPE, reemplaza ejercicios cuando una máquina está ocupada o pospone el día.
    \item Desde la pestaña \textit{Músculos} selecciona un grupo muscular, consulta sus ejercicios y revisa variantes para mantener la continuidad del entrenamiento.
    \item En la pestaña \textit{Rutinas} crea, edita, activa, duplica o archiva sus rutinas.
    \item En la pestaña \textit{Historial} revisa las sesiones realizadas, pospuestas y no realizadas, con el detalle de las series registradas.
\end{enumerate}

Este flujo permite que el usuario continúe su entrenamiento incluso cuando ciertas máquinas o equipos del gimnasio no estén disponibles, ofreciendo alternativas para trabajar el mismo grupo muscular. El siguiente esquema resume la transición entre las zonas principales del sistema.

\begin{figure}[H]
\centering
\footnotesize
\setlength{\fboxsep}{6pt}
\fbox{Inicio} $\rightarrow$ \fbox{Acceso / Registro} $\rightarrow$ \fbox{Onboarding} $\rightarrow$ \fbox{Hoy (sesión)} \\[8pt]
\fbox{Hoy (sesión)} $\leftrightarrow$ \fbox{Rutinas} $\leftrightarrow$ \fbox{Historial} $\leftrightarrow$ \fbox{Músculos} \\[8pt]
\fbox{Músculos} $\rightarrow$ \fbox{Ejercicios por músculo} $\rightarrow$ \fbox{Variantes} \\[8pt]
\fbox{Rutinas} $\rightarrow$ \fbox{Crear / Editar rutina} \qquad \fbox{Perfil} $\rightarrow$ \fbox{Cerrar sesión}
\caption{Esquema del flujo de navegación entre las zonas principales del sistema.}
\label{fig:flujo-navegacion}
\end{figure}

\subsection{Modelo de datos del sistema}

El modelo de datos del sistema organiza la información utilizada por la aplicación. La persistencia se realiza sobre \textit{Cloud Firestore}, donde los datos se dividen en dos ámbitos: un \textbf{catálogo global} compartido por todos los usuarios (músculos, ejercicios, variantes y rutinas prediseñadas) y la \textbf{información personal} de cada usuario, almacenada bajo su identificador (\texttt{users/\{uid\}}). La autenticación y la identidad del usuario son gestionadas por \textit{Firebase Authentication}.

\subsubsection{Entidades principales}

Las entidades del sistema y su contenido se resumen en la Tabla \ref{tab:entidades}.

\begin{table}[H]
\centering
\rowcolors{2}{gymblueSoft}{gymblueSoftAlt}
\arrayrulecolor{gymblue}
\begin{tabularx}{0.98\textwidth}{L{3.4cm} L{3.6cm} Y}
\toprule
\gymtableheaderrow
\gymth{Entidad} & \gymth{Ubicación (colección)} & \gymth{Atributos principales} \\
\midrule
Usuario & Firebase Authentication & Identificador, nombre, correo, foto y proveedor de acceso (correo o Google). \\
Perfil & \texttt{users/\{uid\}/profile} & Objetivo, nivel, días de entrenamiento, equipamiento y estado del onboarding. \\
Músculo & \texttt{muscles} & Nombre, ícono, color e imagen del grupo muscular. \\
Ejercicio & \texttt{exercises} & Músculo asociado, nombre, descripción y equipamiento. \\
Variante & \texttt{variants} & Ejercicio de origen, nombre, descripción y ejercicio sustituto. \\
Rutina prediseñada & \texttt{routine\_templates} & Nombre, descripción, nivel, objetivo, días, equipamiento, división y plan semanal. \\
Rutina & \texttt{users/\{uid\}/routines} & Nombre, ejercicios, estado, plan semanal, días por semana e inicio del ciclo. \\
Día de entrenamiento & \texttt{users/\{uid\}/training\_days} & Fecha, estado del día, ejercicios y series completados. \\
Sesión de entrenamiento & \texttt{users/\{uid\}/workout\_sessions} & Fecha, estado, ejercicios planificados y resultado de la sesión. \\
Registro de ejercicio & \texttt{users/\{uid\}/exercise\_logs} & Fecha, ejercicio y series con peso, repeticiones y RPE. \\
\bottomrule
\end{tabularx}
\caption{Entidades principales del modelo de datos de GymFlow}
\label{tab:entidades}
\end{table}

\subsubsection{Relaciones entre entidades}

Las principales relaciones del modelo se describen en la Tabla \ref{tab:relaciones}.

\begin{table}[H]
\centering
\rowcolors{2}{gymblueSoft}{gymblueSoftAlt}
\arrayrulecolor{gymblue}
\begin{tabularx}{0.98\textwidth}{L{5.2cm} Y}
\toprule
\gymtableheaderrow
\gymth{Relación} & \gymth{Descripción} \\
\midrule
Músculo --- Ejercicio & Un músculo agrupa varios ejercicios; cada ejercicio pertenece a un músculo. \\
Ejercicio --- Variante & Un ejercicio puede tener varias variantes o alternativas para trabajar el mismo músculo. \\
Usuario --- Perfil & Cada usuario posee un único perfil con sus preferencias de entrenamiento. \\
Usuario --- Rutina & Un usuario puede tener varias rutinas, de las cuales como máximo una está activa. \\
Rutina --- Ejercicio & Una rutina referencia un conjunto de ejercicios distribuidos en un plan semanal. \\
Usuario --- Sesión / Registro & Cada usuario genera días y sesiones de entrenamiento, con registros de series por ejercicio y fecha. \\
\bottomrule
\end{tabularx}
\caption{Relaciones entre las entidades del modelo de datos}
\label{tab:relaciones}
\end{table}

Esta organización centraliza la información del catálogo y mantiene aislada la información personal de cada usuario bajo su propio identificador, lo que facilita la seguridad de los datos, la escalabilidad del sistema y la consistencia entre las rutinas, las sesiones y el historial de entrenamiento.
